\documentclass[main]{subfiles}


\title[Evaluation and Benchmarking]{Evaluation and Benchmarking}
\subtitle{Statistical Tests}
\author[Marcel Wever]{Marcel Wever}
\date{}

\titlebackground*{images/AutoML_ML_Evaluation}

%\institute{}
%\date{}

\begin{document}
	
	\maketitle
	

%----------------------------------------------------------------------
%----------------------------------------------------------------------
\begin{frame}{Setting the Scene}
\begin{itemize}
    \item Recall: $\widehat{\riskg}$ is a random variable
    \item Suppose we compare two configurations $\conf^A, \conf^B \in \pcs$ on the same dataset $\dataset$
    \item We observe estimated risks $\widehat{\riskg}^A$ and $\widehat{\riskg}^B$
    \item We want to know: does $\widehat{\riskg}^A \neq \widehat{\riskg}^B$ reflect a true difference?
\end{itemize}
\vspace{.4cm}
\begin{itemize}
    \item Let $d_j = \widehat{\riskg}_j^A - \widehat{\riskg}_j^B$ denote the performance difference on evaluation $j$.\\
    \item \alert{Goal of statistical testing}: Decide based on observed differences $d_1 \ldots, d_k$, whether $\expect[d_j] \neq 0$ or whether the observed differences are only random variation.
\end{itemize}
\end{frame}

\begin{frame}{Why Statistical Testing? A Motivating Example}

\begin{itemize}
    \item Suppose we evaluate $\conf^A$ and $\conf^B$ on $k=10$ CV folds
    \item Observed mean difference: $\bar{d} = \widehat{\riskg}^A - \widehat{\riskg}^B = 0.02$
    \item Is this a real improvement, or explainable by the randomness of the splits?
\end{itemize}

\begin{itemize}
    \item Consider two scenarios with the same $\bar{d}=0.02$:
\end{itemize}

\begin{center}
\begin{tabular}{lll}
     \toprule
     Scenario & Differences $d_j$ & Conclusion\\
     \midrule
     Low variance & all $d_j \approx 0.02$ & Likely a real effect\\
     High variance & $d_j \in \left[ -0.15, 0.19 \right]$ & Likely noise \\
     \bottomrule
\end{tabular}
\end{center}

\begin{itemize}
    \item A raw difference in mean performance is not sufficient to draw conclusions
    \item Must account for the variance of the observed differences
\end{itemize}

\end{frame}

\begin{frame}[c]{Hypothesis Testing Basics}
    \begin{itemize}
        \item Null hypothesis $H_0$: No difference between models (e.g., equal mean performance)
        \item Alternative hypothesis $H_1$: A true difference exists
        \item Significance level $\alpha$: The maximum acceptable probability of incorrectly rejecting $H_0$
        \begin{itemize}
            \item $\alpha$ must be fixed \alert{before} running the test
            \item Commonly $\alpha=0.05$
        \end{itemize}
        \item p-value: Probability of observing results at least as extreme, assuming $H_0$ is true
    \end{itemize}
    \begin{itemize}
        \item Reject $H_0$ is $p \leq \alpha$, where $\alpha$ is the significance level (commonly $0.05$)
        \item \alert{Note}:
    \begin{itemize}
        \item Failing to reject $H_0$ does not prove the two methods are equal
        \item Rejecting $H_0$ does not mean that one is better than the other
    \end{itemize}
    \end{itemize}
\end{frame}

\begin{frame}{Type I and Type II Errors}
    \begin{itemize}
        \item Every hypothesis test can make two kinds of mistakes:
    \end{itemize}
    \begin{center}
        \begin{tabular}{lll}
            \toprule
             & $H_0$ true & $H_0$ false \\
             \midrule
             Reject $H_0$ & \alert{Type I error} (false positive) & Correct (true positive) \\
             Fail to reject & Correct (true negative) & \alert{Type II error} (false negative)\\
             \bottomrule
        \end{tabular}
    \end{center}
    \begin{itemize}
        \item There is a tradeoff: lowering $\alpha$ reduces Type I but increases Type II errors
        \item More data and lower variance in $d_j$ increase power
    \end{itemize}
\end{frame}

\begin{frame}{The $t$-Test}
    \begin{itemize}
        \item Parametric test: assumes differences $d_j$ are i.i.d.\ and normally distributed
        \item \textbf{Paired $t$-Test}: Compare two algorithms on the same $k$ evaluation folds:
        \[
        t = \frac{\bar{d}}{\stddev_d / \sqrt{k}}, \quad \bar{d} = \frac{1}{k} \sum_{j=1}^k d_j, \qquad \stddev_d = \sqrt{\frac{1}{k-1}\sum_{j=1}^k (d_j - \bar{d})^2}
        \]
        \item Null hypothesis: $\mathbb{E}[\bar{d}] = 0$
        \item Sensitive to violations of the normality and the i.i.d.\ assumption
        \item CV folds are not independent (sharing training data)
        \item Corrected resampled $t$-test \liturl{Nadeau \& Bengio 2003}{https://link.springer.com/article/10.1023/A:1024068626366}: Adjusts for the dependency of CV folds to avoid inflated type-I error
    \end{itemize}
\end{frame}

\begin{frame}{Non-Parametric Tests}
    \begin{itemize}
        \item Do not assume a specific distribution of the test statistics 
        \item Still assume that observations are i.i.d.
        \item More robust than parametric tests, but generally \alert{less powerful}: require larger or more consistent differences to reject $H_0$
        \item \textbf{Wilcoxon Signed-Rank Test:} Paired test comparing two algorithms across datasets
        \begin{itemize}
            \item Ranks the absolute differences and accounts for their sign
            \item Preferred over paired $t$-test when normality cannot be assumed
        \end{itemize}
        \item \textbf{Sign test:} Simplest paired test, only considers the direction of differences
        \begin{itemize}
            \item Less powerful than Wilcoxon, but very robust to outliers
        \end{itemize}
    \end{itemize}
\end{frame}

\begin{frame}{Ombnibus Tests and Multiple Comparisons}
    \begin{itemize}
        \item With $M>2$ configurations, running all $\binom{M}{2}$ pairwise tests inflates the Type I error:
        $$P(\text{at least one false positive}) = 1 - (1-\alpha)^{\binom{M}{2}}$$
        \item First, apply an \alert{omnibus test} to check for \alert{any} difference among all $M$ configurations
        \item \textbf{Friedman test}: Omnibus test across $M$ configurations on $N$ datasets or splits
        \begin{itemize}
            \item Ranks configurations per dataset, tests whether mean ranks differ across all $M$
            \item Rejects $H_0$: ``all configurations perform equally'' without specifying which pair differs
        \end{itemize}
        \item If Friedman rejects $H_0$, follow up with \alert{post-hoc} pairwise tests
    \end{itemize}
\end{frame}

\begin{frame}{Post-Hoc Tests and Corrections}
    \begin{itemize}
        \item Post-hoc tests perform pairwise comparisons after an omnibus test, with correction for multiple testing
        \item \textbf{Bonferroni}: Divide $\alpha$ by the number of tests. Simple but increases Type II error
        \item \textbf{Holm}: Step-down procedure: sort p-values, apply decreasing thresholds
        \item \textbf{Nemenyi test}: All pairwise comparisons after Friedman; based on rank differences; convenient but less powerful than Holm
    \end{itemize}
    
    \vspace{.2cm}
    
    \begin{minipage}{.6\textwidth}
    \begin{itemize}
        \item Results of Friedman + post-hoc often visualized via \alert{critical difference} (CD) diagrams
    \end{itemize}
    \end{minipage}
    \hfill
    \begin{minipage}{.3\textwidth}
    \resizebox{\textwidth}{!}{
        \usetikzlibrary{arrows.meta}

\begin{tikzpicture}[
  >=Stealth,
  font=\small,
  pin line/.style  = {draw=black, line width=0.6pt},
  sig bar/.style   = {draw=black, line width=4pt, line cap=butt},
  axis line/.style = {draw=black, line width=1.2pt},
  tick/.style      = {draw=black, line width=0.8pt},
]

% ── Coordinate system ──────────────────────────────────────────────────────
% Ranks 1..5 mapped linearly onto x = 0..9 cm
\def\axisW{9}
\def\kmax{5}
\pgfmathsetmacro{\scaleX}{\axisW / (\kmax - 1)}  % 2.25 cm per rank unit

\def\axisY{0}

% x positions of each method
\pgfmathsetmacro{\xA}{(1.500 - 1) * \scaleX}   % Method 1
\pgfmathsetmacro{\xB}{(2.300 - 1) * \scaleX}   % Method 2
\pgfmathsetmacro{\xC}{(3.033 - 1) * \scaleX}   % Method 3
\pgfmathsetmacro{\xD}{(3.967 - 1) * \scaleX}   % Method 4
\pgfmathsetmacro{\xE}{(4.200 - 1) * \scaleX}   % Method 5

% CD bracket half-width
\pgfmathsetmacro{\CDhalf}{1.1137 / 2 * \scaleX}

% ── Significance bars ──────────────────────────────────────────────────────
% Non-sig pairs (diff < CD = 1.1137):
%   Method 1--2 (diff=0.800), Method 2--3 (diff=0.733), Method 3--4 (diff=0.933)
%
% Bar slot assignment (bars hang below axis):
%   Bar 1 [M1--M2]: slot y = -0.7
%   Bar 2 [M2--M3]: overlaps Bar 1 horizontally -> slot y = -1.25
%   Bar 3 [M3--M4]: no overlap with Bar 1 -> slot y = -0.7

\def\barYone{-0.7}
\def\barYtwo{-1.25}

% Bar 1: Method 1 -- Method 2
\draw[sig bar] (\xA, \barYone) -- (\xB, \barYone);
% Bar 2: Method 2 -- Method 3
\draw[sig bar] (\xB, \barYtwo) -- (\xC, \barYtwo);
% Bar 3: Method 3 -- Method 4
\draw[sig bar] (\xC, \barYone) -- (\xD, \barYone);

% ── Label rows ─────────────────────────────────────────────────────────────
% Each method gets its own row, staggered to avoid overlap.
% Rows are assigned by inspecting x-distances: methods closer than ~1.8cm
% need to be on different rows.
%
% x positions (approx): M1=0.00, M2=2.92, M3=4.57, M4=6.68, M5=7.20
% M4 and M5 are only 0.52 cm apart -> different rows
% All others are >1.8 cm apart -> can share a row
%
% Row assignment (row 1 = highest, i.e. closest to axis):
%   Row 1 (y = -1.85): Method 1, Method 2, Method 3, Method 4
%   Row 2 (y = -2.45): Method 5

\def\rowOne{-1.85}
\def\rowTwo{-2.45}

% Pin bottom per method: reach below the lowest bar they belong to,
% then continue to their label row.
%   Method 1: Bar1 (y=-0.70) -> pin to row1
%   Method 2: Bar1+Bar2 (y=-1.25) -> pin to row1
%   Method 3: Bar2+Bar3 (y=-1.25) -> pin to row1
%   Method 4: Bar3 (y=-0.70) -> pin to row1
%   Method 5: no bar -> pin to row2

% Method 1
\filldraw[black] (\xA, \axisY) circle (2.2pt);
\draw[pin line]  (\xA, \axisY) -- (\xA, \rowOne);
\node[below, font=\small] at (\xA, \rowOne) {Method 1};

% Method 2
\filldraw[black] (\xB, \axisY) circle (2.2pt);
\draw[pin line]  (\xB, \axisY) -- (\xB, \rowOne);
\node[below, font=\small] at (\xB, \rowOne) {Method 2};

% Method 3
\filldraw[black] (\xC, \axisY) circle (2.2pt);
\draw[pin line]  (\xC, \axisY) -- (\xC, \rowOne);
\node[below, font=\small] at (\xC, \rowOne) {Method 3};

% Method 4
\filldraw[black] (\xD, \axisY) circle (2.2pt);
\draw[pin line]  (\xD, \axisY) -- (\xD, \rowOne);
\node[below, font=\small] at (\xD, \rowOne) {Method 4};

% Method 5  (bumped to row 2 to avoid overlap with Method 4)
\filldraw[black] (\xE, \axisY) circle (2.2pt);
\draw[pin line]  (\xE, \axisY) -- (\xE, \rowTwo);
\node[below, font=\small] at (\xE, \rowTwo) {Method 5};

% ── Rank axis ──────────────────────────────────────────────────────────────
\draw[axis line] (0, \axisY) -- (\axisW, \axisY);

\foreach \r in {1,...,5} {
  \pgfmathsetmacro{\tx}{(\r - 1) * \scaleX}
  \draw[tick] (\tx, \axisY-0.12) -- (\tx, \axisY+0.12);
  \node[above, yshift=2pt] at (\tx, \axisY+0.12) {\r};
}

\node[above, yshift=10pt] at (0.5*\axisW, \axisY+0.12)
  {\textit{Average Rank \enspace (1 = best)}};

% ── CD bracket ─────────────────────────────────────────────────────────────
\pgfmathsetmacro{\cdCX}{\axisW - \CDhalf}
\def\cdY{0.9}

\draw[{Stealth[length=5pt]}-{Stealth[length=5pt]}]
  (\cdCX-\CDhalf, \cdY) -- (\cdCX+\CDhalf, \cdY);
\draw[tick] (\cdCX-\CDhalf, \cdY-0.1) -- (\cdCX-\CDhalf, \cdY+0.1);
\draw[tick] (\cdCX+\CDhalf, \cdY-0.1) -- (\cdCX+\CDhalf, \cdY+0.1);
\node[above] at (\cdCX, \cdY) {$\mathrm{CD} = 1.114$};

% ── Caption ────────────────────────────────────────────────────────────────
% \node[below, align=center, font=\footnotesize] at (0.5*\axisW, \rowTwo-0.55)
  % {Connected methods are not significantly different
   % (Nemenyi, $\alpha=0.05$, $\chi^2_{F}=61.39$, $p=1.48\times10^{-12}$, $N=30$).};

\end{tikzpicture}
}
    \end{minipage}
\end{frame}

\begin{frame}{Practical Recommendations}
\alert{Choosing the right test}:
\begin{itemize}
    \item Two configurations, paired results: Wilcoxon signed-rank test (or corrected $t$-test)
    \item More than two configurations: Friedman omnibus test + post-hoc with correction (Holm or Nemenyi)
    \item Always report effect sizes alongside p-values where possible
\end{itemize}
\textcolor{red}{Common pitfalls}:
\begin{itemize}
    \item Using an uncorrected $t$-test on CV folds (inflated type-I error)
    \item Fixing $\alpha$ after seeing the data
    \item Interpreting $p > \alpha$ as evidence that configurations are equivalent
    \item Running many pairwise tests without correction: accumulates false positives
    \item Choosing the test after seeing the data
\end{itemize}
\end{frame}

\begin{frame}{Summary} 
Most important insights from this video:

\begin{enumerate}
    \item Performance differences are random variables; statistical tests determine whether observed differences are consistent with chance
    \item Type I errors are controlled by $\alpha$; Type II errors are reduced by increasing power through more data or lower variance
    \item Wilcoxon Signed-Rank Test is a robust default for comparing two configurations
    \item For more than two configurations, always use an omnibus test (Friedman) first, followed by post-hoc tests with multiple testing correction
    \item Never run pairwise tests without correction, and never fix $\alpha$ after seeing the results
\end{enumerate}

\end{frame}

\end{document}
